\documentclass[11pt]{article}
\usepackage{../common/indexus}
\usepackage{../figkit/ixfig}

\begin{document}

\title{Distributed Prefix Index Protocol (DPIP)\\[3pt]{\large\mdseries Multiscale Aggregation over a Peer-to-Peer Index}}
\author{\ixbyline}
\hypersetup{pdftitle={Distributed Prefix Index Protocol (DPIP). Multiscale Aggregation over a Peer-to-Peer Index}, pdfauthor={Rémi Henriot (AI Co-authored)}}
\date{November 2026}
\maketitle

\begin{abstract}
\noindent \textbf{Abstract. }In DPIP, peers share responsibility for an index. Each location prefix names a zone within a collection. Extending that prefix lets a reader move from a regional total to a smaller area. A house, for example, has one location in a tree. A reader can ask for the street listing or for the number of houses in its district. We follow one listing from its coordinates to a durable acknowledgement, through a change of host, and back to a map. A tree address gives the reader a destination to try. A right recorded on disk determines which node may accept the write. Each parent combines versioned contributions from its children, so a repeated message cannot add the house twice.
\vspace{1em}\end{abstract}

\section{Introduction}
Data becomes useful when it can be found: the houses for sale near an address, the weather at a given hour, the news of the day scrolled back in time. Finding data requires an index. Building one is costly, so in practice one party builds it and the others query it. That party then decides what is found, in which order and on which terms.

Indexus is the project that originated the DPIP specification and developed its reference implementation. Any party can run part of it, any reader can compute the address of a zone, and readers following the same complete data versions compute the same answer. For maps, feeds and dashboards, the index also answers totals, how many and how much in a region, and those totals compose from versioned contributions. A summary may still lag an acknowledged write: correctness of its calculation and freshness of its inputs are separate.

Henriot's 2024 papers introduced the tree, the encoding of locations and the incremental search~[\ixref{1},~\ixref{2}]. This paper describes the protocol as it is built today.

\paragraph{One example throughout.}
A producer adds one house to collection BB. Each listing has a position and a price. A city view asks for counts and price totals in large cells; a street view asks for smaller cells. The records stay in one tree. The reader changes the depth at which it reads, rather than building another index for each zoom.

\paragraph{Companion papers.}
\ixpaper{client} explains map refinement and nearby search.
\ixpaper{network} details \ixsection{network}{handoff}{installed authority and granular transfer}.
Definitions and workloads behind the figures are in
\ixsection{properties}{benchmarks}{the properties paper's benchmark section};
\ixpaper{economics} describes proposals, and \ixpaper{history} records provenance.

\section{Tree}\hypertarget{ix:tree}{}
DPIP places the network's nodes, collections, zones and records in one binary tree.

\begin{figure}[H]
    \centering
    \ixfig[10cm]{../figkit/orig/tree.tikz}
    \caption{Binary tree}
\end{figure}

We write positions with characters. In our examples a character is 2 bits, written A, B, C or D. In practice a character is 6 bits, from an alphabet of 64 symbols, and a full address is 96 bits, 16 characters.

\begin{figure}[H]
    \centering
    \ixfig[10cm]{../figkit/orig/encoded.tikz}
    \caption{Encoded tree (BASE4)}
\end{figure}

A prefix names a whole branch: C includes CA, CB, CC and CD in our four-letter drawings. A deeper prefix names a smaller part of the same tree. The drawings use four letters so that all branches fit on a page; the implementation uses 64.

\section{Space}\hypertarget{ix:space}{}
A record has coordinates: a latitude and a longitude, a time, or any measurable value. A \textbf{space} turns coordinates into a position in the tree by halving its dimensions one after the other. Each halving gives one bit: 1 above the midpoint, 0 at or below it. The \textbf{mask} says which dimension each bit cuts. At each level the \textbf{offset} adds the corresponding collection-name digit, modulo the alphabet size. Two collections can therefore give the same coordinates different location paths.

\begin{figure}[H]
    \centering
    \ixfig[12cm]{../figkit/orig/masks.tikz}
    \caption{Two subdivisions of the same space}
\end{figure}

Our house lies at latitude $-20$ and longitude $45$. Starting with the full globe, the first latitude cut puts it in the lower half, and the first longitude cut in the upper half. These bits give $01$. In collection BB the first offset is B, or one. The sum in binary is $01+1=10$, written C. We store the listing at this precision, C. For a finer reading, the next two cuts compare the same coordinates with $-45$ and $90$. They give $10$; adding the next B gives $11$, or D. At that precision the location would be CD, a child of C.

\section{Collection}\hypertarget{ix:collection}{}
A \textbf{record} is one entry: an identifier, a location and a few numbers, such as a count or a price. A \textbf{collection} is a named set of records that use the same space. A \textbf{zone} is a part of a collection: all the records whose location starts with the same characters. The zone C holds every record whose location starts with C, and the zone CB every record starting with CB. A map zoomed out reads large zones; zoomed in, it reads small ones.

Each zone keeps its \textbf{summary}: a count and one aggregate per numeric column. An ordinary listing contributes one to the count. A pre-aggregated row may carry a weight, so one stored row can represent several observations. By default each number is summed; a collection may also declare \texttt{min}, \texttt{max}, \texttt{first}, \texttt{last}, or linear-recurrence pairs \texttt{affine} $(a, b)$ combined in cell order. With associative operations and a consistent set of child versions, folding a complete cut in cell order gives the root in exact arithmetic. Floating-point sums may round differently when regrouped. For a self-certifying collection, a canonical \texttt{IXC1} description fixes the geometry, bounds, bit schedule and column operations. A reader verifies the description by deriving its 96-bit name. A collection named by hand, as in the Anglet sample, needs its schema from the producer.

\begin{figure}[H]
    \centering
    \ixfig[15cm]{../figkit/orig/collection.tikz}
    \caption{Collection, set and items}
\end{figure}

The 2024 papers called zones \emph{sets} and records \emph{items}. Before our house is added, collection BB holds 60 records. Zones A, B, C and D account for 15, 3, 33 and 9 respectively. A map can draw these four counts without reading the records.

\section{Network}
A host can run several virtual nodes, or \textbf{vnodes}. Each vnode has its own address, a leaf of the tree. Each node in the drawings represents one such participant.

\begin{figure}[H]
    \centering
    \ixfig[10cm]{../figkit/orig/p2p-network.tikz}
    \caption{Network of nodes}
\end{figure}

A zone is identified by its collection and its location path. For routing, the proposed rule replaces the leading digits of the collection name with that path and keeps the unused suffix. The root, written @, uses no location digits. Here is what happens as one follows the BB collection down two branches.

\begin{center}
\begin{tabular}{@{}llcl@{}}
\toprule
Path & Location digits & Collection suffix & Routing address \\
\midrule
@  & none & BB & BB \\
C  & C    & B  & CB \\
CB & CB   & none & CB \\
CD & CD   & none & CD \\
\bottomrule
\end{tabular}
\end{center}

C and its child CB lead to the same routing address. The B that the parent borrowed from the collection is exactly the digit its child supplies. Yet a count for C includes CA, CB, CC and CD, while a count for CB covers only that branch. Our listing stored at C contributes to C and @. Of the nodes in the drawing, CA has the smallest bitwise exclusive-or (XOR) distance to CB. A reader may try CA for either path and must still specify which zone it wants.

\begin{figure}[H]
    \centering
    \ixnew{figures/owner.tikz}
    \caption{Zone C of BB has address CB. CA is the nearest node shown and holds its write permission in this example.}
\end{figure}

The address predicts a destination. A write succeeds there when the disk holds the installed right for the requested zone. The \ixsection{network}{authority}{network paper} follows the case where a child has moved and its parent is still served by the donor. It also records a byte-packing difference between the proposed address rule and the current codec.

\section{Routing}
Nodes exchange membership updates in the background. An entry also learns from a successful write which node applied it. This helps when the holder of our listing changes before all the readers have heard about the new vnode. The next request can use the learned route. The destination checks its own permission, so a write can complete while membership updates are still in flight.

The 2024 drawing groups contacts by the first bit where their paths differ, following the XOR geometry also used in Kademlia~[\ixref{3}]. It illustrates how to choose contacts at several distances.

\begin{figure}[H]
    \centering
    \ixfig[13cm]{../figkit/orig/p2p-routing.tikz}
    \caption{Routing contacts grouped by the first differing bit}
\end{figure}

\Needspace{10\baselineskip}
\section{Writing}
For the new house the write follows four steps.
\begin{enumerate}[leftmargin=1.5em,itemsep=1pt]
  \item The client sends the record to any node, with an identifier for the request.
  \item If that node does not hold the zone, it forwards the record to the node that does.
  \item The holder stores the record on its disk, together with a receipt of the request.
  \item The holder answers. This answer is the \textbf{acknowledgement}.
\end{enumerate}

An acknowledged record is on disk. If the client gets no answer, it sends the same request again. The holder recognises the identifier and returns the first answer: the record is never added twice.

Suppose the producer already has a town's listings before the first reader arrives. It can sort the records offline into checksummed parts (\texttt{IXP1}). In the default fresh-import mode, the target zone must be empty on the receiving disk, and the receiver must hold the right to write each addressed zone. A concurrent hand-over can cause the import to be refused. One local installation publishes the whole batch with one receipt; retrying the same batch leaves it in place. A price change tomorrow uses the ordinary write path. The \href{https://github.com/indexus/core/blob/5188b61/internal/storage/bulk_test.go}{storage tests} exercise both replay and crash recovery.

\section{Delegation}
These three drawings show how a collection spread through the network in the 2024 example. Thick lines lead to nodes, dashed lines lead to zone addresses, and darker tiles indicate zones in use.

\begin{figure}[H]
    \centering
    \ixfig[13cm]{../figkit/orig/p2p-delegation-1.tikz}
    \caption{First, BC serves the root of collection BB.}
\end{figure}

\begin{figure}[H]
    \centering
    \ixfig[13cm]{../figkit/orig/p2p-delegation-2.tikz}
    \caption{Next, AB serves A and CA serves C. BC still serves the root.}
\end{figure}

\begin{figure}[H]
    \centering
    \ixfig[13cm]{../figkit/orig/p2p-delegation-3.tikz}
    \caption{Later, CA still serves C, while CC serves the smaller zones CC and CD.}
\end{figure}

A busy vnode can delegate part of its address space to a new vnode. Which host runs it determines where the load lands. Moving work between two vnodes on one machine can leave that machine just as busy. A transfer proceeds in lots, with descendants first and hotter zones preferred at the same depth. The new vnode can serve each installed lot while the donor handles the rest.

\section{Moving a zone}\hypertarget{ix:handoff}{}
When a zone changes holder, the two nodes may for a short time disagree on who holds it. If both accepted writes, the zone would exist in two versions. So the right to write is handed over explicitly:
\begin{enumerate}[leftmargin=1.5em,itemsep=1pt]
  \item The receiver reserves a ticket. The holder copies the zone and catches up recent writes while still accepting more.
  \item The holder fences this lot on disk. It sends the last changes and the seal for that version.
  \item The holder durably authorises the receiver under the same ticket.
  \item The receiver durably installs the data and permission together. It can now accept writes for this zone.
\end{enumerate}

\begin{figure}[H]
    \centering
    \ixnew{figures/handoff.tikz}
    \caption{A zone moves. The donor fences before the receiver installs. The interval between them has no writer.}
\end{figure}

During an ordinary lot handover the pause concerns the zones in that lot. For example, after lot A is installed, its writes go to the new node while lots B, C and D still use the old one. Waiting for all four before enabling A would throw away the benefit of granular transfer.

\begin{figure}[H]
    \centering
    \ixnew{figures/progressive.tikz}
    \caption{Progressive relief: writes follow each installed and activated lot. Four lots illustrate the mechanism; equal-sized cells do not imply equal traffic or equal transfer times.}
\end{figure}

A smaller lot limits the work in one hand-over; it does not alone bound its duration. Disk latency, the writes to catch up and the network still matter. The implementation can also fence the residual region to finish a join under continuous zone creation. Thus early service and complete join are different milestones, and a universal freeze bound is not yet established.

Receipts move with the zone, so a request sent again after the move is still recognised. Each step is written on disk before the next one starts, and is resumed after a restart.

\section{Sums across nodes}\hypertarget{ix:summaries}{}
Our new listing raises the count of C from 33 to 34. The vnode that holds C sends its absolute count and version to the holder of its parent. When the parent receives the newer version, it replaces 33 with 34, taking the root from 60 to 61. A duplicate or a late copy of the old version leaves the total unchanged.

\begin{figure}[H]
    \centering
    \ixnew{figures/summary.tikz}
    \caption{Zone C, held by CA, sends its sum to zone @, held by BC.}
\end{figure}

A node can batch several changes in one summary. A map may therefore still show 60 after our house was acknowledged. A refreshed normal read must cover the relevant holders and receivers, or report that it cannot. The \ixsection{client}{freshness}{client paper} explains when a map refresh actually reaches the network.

\section{Search}\hypertarget{ix:search}{}
A map reads the zones of the area it shows, at the size given by its zoom. A feed, or any search around a point, reads zones in order of distance from its origin, layer by layer. The sums of the zones tell the search how many records lie within each distance, so it stops as soon as it has enough records closer than everything not yet read. This is the algorithm of Henriot's 2024 search paper, now used by the reference client.

\begin{figure}[H]
    \centering
    \ixfig[15cm]{../figkit/orig/layers.tikz}
    \caption{Layers and concentric circles}
\end{figure}

Collections that use the same space can be read together: a map can show property sales, the weather and a person's own trips, from different producers, without copying any of them.

\section{What has been measured}
In two AWS runs with direct routing, each 70-second plateau at 16\,000 requested writes/s gave a client write p99 of 16.384\,ms, or 32.768\,ms measured from the scheduled arrival. The \ixsection{properties}{benchmarks}{benchmark paper} documents the workloads, sampled rereads and limits of the audits.

\section{Conclusion}
The producer writes one listing. Its address tells an entry where to try, its holder records the listing and a receipt, and a later transfer carries both. The city's count increases when its parent receives the new version for C. These steps allow the next map to reuse the index built by the first. A durable local acknowledgement covers restart with the same disk. Protection against permanent disk loss needs a separately qualified replica. Verification of a total from an untrusted host remains a proposal.

\Needspace{5\baselineskip}
\paragraph{Contributions.} Rémi Henriot introduced the tree, encoding and search in his 2024 papers. AI-assisted drafting of this edition follows one house through those mechanisms and relates the example to the measured hand-off. The \ixsection{history}{ai}{project history} dates the contributions of the tools used.

\section*{References}
\begin{enumerate}[label={[\arabic*]},leftmargin=2em,itemsep=1pt]
  \item \hypertarget{reference:1}{}R. Henriot. \ixbib{search2024}. 2024.
  \item \hypertarget{reference:2}{}R. Henriot. \ixbib{network2024} (first version). 2024.
  \item \hypertarget{reference:3}{}P. Maymounkov and D. Mazi\`eres. Kademlia: A Peer-to-peer Information System Based on the XOR Metric. 2002.
\end{enumerate}

\end{document}
